\chapter{Instalación paso a paso}
\label{cap:instalacion}

\begin{enpocas}
Qué hay que instalar y en qué orden: ROS~2 y Gazebo, PX4 con el módulo neuronal, el agente uXRCE-DDS,
el workspace, ArduPilot y los entornos de Python. Todo está en \cod{scripts/instalar_entorno.sh},
pensado para correrse bloque por bloque y sin \cod{sudo} salvo donde se indica.
\end{enpocas}

\section{0. ROS 2 Jazzy y Gazebo Harmonic}

Ubuntu 24.04 con ROS~2 Jazzy y el puente \cod{ros-jazzy-ros-gz} (que trae Gazebo Harmonic), instalados
con el repositorio oficial de ROS~2 (\cod{ros-jazzy-desktop}, \cod{ros-dev-tools},
\cod{ros-jazzy-ros-gz}).

\section{1. PX4 release/1.17}

\begin{lstlisting}[style=bash]
cd ~ && git clone --recursive -b release/1.17 https://github.com/PX4/PX4-Autopilot.git
python3 -m venv --system-site-packages ~/.venvs/px4
~/.venvs/px4/bin/pip install -r ~/PX4-Autopilot/Tools/setup/requirements.txt ninja
\end{lstlisting}

\begin{aviso}
La etiqueta \cod{v1.17.0} de PX4 no existe (solo \cod{v1.17.0-rc2}): se usa la rama \cod{release/1.17}.
Ubuntu 24.04 no deja instalar con \cod{pip install --user} sobre el Python del sistema (PEP 668): por eso
cada grupo de dependencias va en su propio entorno.
\end{aviso}

\section{2. Agente Micro XRCE-DDS 2.4.3}

\begin{lstlisting}[style=bash]
cd ~ && git clone -b v2.4.3 https://github.com/eProsima/Micro-XRCE-DDS-Agent.git
cd Micro-XRCE-DDS-Agent && mkdir build && cd build
env -i HOME=$HOME PATH=/usr/bin:/bin cmake .. -DCMAKE_INSTALL_PREFIX=$HOME/.local
env -i HOME=$HOME PATH=/usr/bin:/bin make -j"$(nproc)" && make install
\end{lstlisting}
Se compila con un entorno limpio (\cod{env -i}) para que use su propio Fast-DDS y no el de ROS, y se
instala en \cod{~/.local} sin \cod{sudo}.

\section{3. Workspace}

\begin{lstlisting}[style=bash]
git clone https://github.com/jairoroke/<repositorio>.git ~/ros2_ws
cd ~/ros2_ws/src
git clone -b release/1.17 https://github.com/PX4/px4_msgs.git
git clone https://github.com/ArduPilot/ardupilot_gazebo.git       # solo para piloto:=ardupilot
cd ~/ros2_ws && colcon build --symlink-install && source install/setup.bash
\end{lstlisting}

\section{4. Entornos de Python}

\begin{lstlisting}[style=bash]
python3 -m venv ~/venv-mision
~/venv-mision/bin/pip install numpy scipy matplotlib pillow tensorflow-cpu    # exportar la red
\end{lstlisting}
\begin{itemize}
  \item \cod{~/venv-mision}: TensorFlow, para exportar la red a TFLite.
  \item \cod{~/.venvs/yolo}: Ultralytics con PyTorch y CUDA (y \cod{numpy<2}, para leer los mensajes de
        ROS~2). Corre el detector y el refinamiento PPO.
  \item \cod{~/venv-ardupilot}: MAVProxy, pymavlink y \cod{python-xlib} (para acomodar las ventanas).
\end{itemize}

\section{5. Vehículo, red y compilación de PX4}

\begin{lstlisting}[style=bash]
cd ~/ros2_ws/src/mision
python3 scripts/generar_vehiculo.py      # tarot650_busqueda, x500_busqueda y sus variantes para ArduPilot
cd codigo
python3 nn_imitacion.py --vehiculo tarot650 --masa 2.25 --com 0.00067 0.00067 --adelanto \
    --salida pesos_nn_tarot650_busqueda.npz --figura ../imagenes/nn_imitacion_tarot650_busqueda.pdf
~/venv-mision/bin/python exportar_tflite.py --pesos pesos_nn_tarot650_busqueda.npz --salida control_net_tarot650_busqueda
bash ../scripts/instalar_px4.sh ~/ros2_ws/src/mision/codigo/control_net_tarot650_busqueda
\end{lstlisting}
\cod{instalar_px4.sh} copia a \cod{~/PX4-Autopilot} los vehículos \cod{tarot650_busqueda} y
\cod{x500_busqueda}, los \emph{airframes} \cod{4100_gz_tarot650_busqueda} y \cod{4099_gz_x500_busqueda}
(carpeta \cod{px4/}) y la red (guardando una vez la original como \cod{control_net.*.px4}), agrega
\cod{/fmu/out/neural_control} a los tópicos que exporta PX4 y compila \cod{make px4_sitl_neural}.
Repite estos pasos cada vez que cambies la red o el vehículo.

\begin{aviso}
PX4 lleva \textbf{una sola red compilada}: ahora es la del Tarot. Para volar el x500 con
\cod{controlador:=red} en PX4 hay que reinstalar \cod{control_net_x500_busqueda}. Con ArduPilot no
hace falta: \cod{controlador_red} carga los pesos del vehículo que se lanza.
\end{aviso}

\section{6. ArduPilot (opcional)}

ArduCopter 4.7.2 compilado para SITL en \cod{~/ardupilot} (\cod{./waf configure --board sitl \&\& ./waf
copter}) y el plugin \cod{ardupilot_gazebo} en el workspace.

\section{Comprobación}

\begin{lstlisting}[style=bash]
cd ~/ros2_ws/src/mision && python3 -m pytest -q test/          # 8 pruebas
ros2 launch mision simulacion.launch.py gui:=false              # una misión completa, unos 4 min
\end{lstlisting}
